\section{有限网络的极限与自相似}
\subsection{先定义截断，再谈无限}

无限电阻网络不是把有限网络的末端随意删掉，而是指定一列逐渐增大的有限网络及其边界条件，再考查指定端口的电压电流关系是否收敛。末端开路、末端短路和周期边界在有限尺寸下给出不同结果；在有足够耗散且影响向远处衰减的梯形网络中，它们可能趋于同一局部极限，但这一点应证明或检查。

求自相似电阻时，设剩余部分的端口电阻为 $X$，切口应落在明确的两个节点上。移去一个单元后，剩余网络必须与原网络有相同的端口定义和电阻比例，才可重新记为 $X$。若从“先串联一段”改成“先并联一个横档”，端口电阻往往相差一个串联电阻。许多相差 $r$ 的答案正是由切口不同造成。

\fig{c4:single-cut}{单有阻导轨的输入切口：$A$ 先经 $r$ 到 $P$，下导线各点属于同一节点 $B$}{\begin{tikzpicture}

\draw(0,2)to[R,l={$r$}](3,2)to[R,l={$r$}](6,2)to[R,l={$r$}](9,2);
\draw(0,0)--(9,0);
\foreach \x in {3,6,9}{\draw(\x,2)to[R,l={$r$}](\x,0);}
\draw[dashed](9,2)--(10,2)(9,0)--(10,0);
\fill(0,2)circle(1.5pt)node[left]{$A$};\fill(0,0)circle(1.5pt)node[left]{$B$};
\fill(3,2)circle(1.5pt)node[above]{$P$};
\draw[guide](3.6,-.6)rectangle(10.4,2.8);\node at(7,3.4){从 $P,B$ 向右的尾部与 $A,B$ 输入相同};

\end{tikzpicture}}
图\ref{c4:single-cut}的虚框包含从 $P$ 向右再经过一只串联电阻的尾部，不含 $P$ 处的竖档。因此在 $P,B$ 之间，本地竖档 $r$ 与尾部输入 $X$ 并联；从 $A$ 输入还须加左侧串联 $r$。相似的几何图形只有在切口和端点也对应时，才能使用同一个递推变量。

\subsection{固定点的物理筛选与收敛}

例如单根有阻上导轨、理想下导线构成的梯形网络，每单元先串联 $r$，再接一个对下导线的横档 $r$。移去首单元后剩余输入电阻仍为 $X$，故
\[
 X=r+\frac{rX}{r+X},\qquad X^2-rX-r^2=0.
\]
正根为 $(1+\sqrt5)r/2$，负根不能表示正电阻网络。仅丢弃负根还不够：应说明有限网络递推 $X_{n+1}=f(X_n)=r+rX_n/(r+X_n)$ 会逼近该根。对 $X$ 求导时 $r$ 为常量，按商的求导法则，
\[
 f'(X)=0+\frac{r(r+X)-rX\cdot1}{(r+X)^2}
 =\frac{r^2}{(r+X)^2}.
\]
在 $X\ge r$ 时有 $0<f'(X)\le1/4$。设正固定点为 $X_*$，中值定理给出
$|X_{n+1}-X_*|=|f'(\xi_n)|\,|X_n-X_*|\le|X_n-X_*|/4$，其中 $\xi_n$ 位于 $X_n$ 与 $X_*$ 之间。第一节以后各值都落入该范围，故误差最多按 $4^{-n}$ 衰减，固定点确实是有限网络的极限。

\fig{c4:truncation-bounds}{相同的串并联递推作用于开路、短路两种末端，有限值从两侧夹逼同一个被动固定点}{\begin{tikzpicture}

\draw[axis](0,0)--(6.8,0)node[right]{$N$};
\draw[axis](0,0)--(0,3.6)node[above]{$X_N/r$};
\draw[blue,thick](1,3)--(2,2.5)--(3,2.4375)--(4,2.428571)--(5,2.427273)--(6,2.427083);
\draw[red,thick](1,1.5)--(2,2.25)--(3,2.4)--(4,2.423077)--(5,2.426471)--(6,2.426966);
\foreach \x/\y in {1/3,2/2.5,3/2.4375,4/2.428571,5/2.427273,6/2.427083}{\fill[blue](\x,\y)circle(1.5pt);}
\foreach \x/\y in {1/1.5,2/2.25,3/2.4,4/2.423077,5/2.426471,6/2.426966}{\fill[red](\x,\y)circle(1.5pt);}
\draw[guide](0,2.427051)--(6.4,2.427051);
\node[left]at(0,2.427051){$\varphi$};\node[left]at(0,1.5){$1$};\node[left]at(0,3){$2$};
\foreach \n in {1,...,6}{\draw(\n,0)--(\n,-.1)node[below]{$\n$};}
\node[blue]at(3.3,3.4){末端开路 $X_0=\infty$};
\node[red]at(3.3,1.3){末端短路 $X_0=0$};

\end{tikzpicture}}
图\ref{c4:truncation-bounds}采用 $X_{N+1}=f(X_N)=r+rX_N/(r+X_N)$。开路末端依次给 $X_1=2r,X_2=5r/3$；短路末端依次给 $X_1=r,X_2=3r/2$。由于 $f$ 单调递增，开路截断给上界、短路截断给下界；两序列进入 $X\ge r$ 后，导数上界 $1/4$ 使两界距离趋零。这里 $\varphi=(1+\sqrt5)/2$。有限两节仍有不同的端口电阻，只有极限可以合一。

若把输入端改在第一横档上，则电阻为 $r\parallel X=X-r=(\sqrt5-1)r/2$。这两个值描述不同端口，不能因为图形都称“无限梯形”就只记其中一个常数。

\typeout{ECFIG:fig:c4:cut}

\section{递推方程与电流的空间衰减}
\subsection{用节点电势建立差分方程}

单上导轨网络中，每相邻节点间电阻 $r$，每节点经横档 $R$ 接公共理想下导线。取下导线为零电势，在无外部注入的第 $n$ 个节点列守恒：
\[
 \frac{V_n-V_{n-1}}r+\frac{V_n-V_{n+1}}r+\frac{V_n}R=0.
\]
\fig{c4:node-difference}{在第 $n$ 个上轨节点，把三条有阻支路均规定为流出方向}{\begin{tikzpicture}

\draw(-3,0)to[R,l={$r$}](0,0)to[R,l={$r$}](3,0);
\draw(0,0)to[R,l={$R$}](0,-2.8);\draw(-3,-2.8)--(3,-2.8);
\fill(0,0)circle(1.5pt)node[above=24pt]{$V_n$};
\node[left]at(-3,0){$V_{n-1}$};\node[right]at(3,0){$V_{n+1}$};
\node[below]at(0,-2.8){公共理想线：$0$};
\draw[flow](-.3,.6)--(-1.4,.6);\draw[flow](.3,.6)--(1.4,.6);
\draw[flow](.6,-.7)--(.6,-1.8);

\end{tikzpicture}}
图\ref{c4:node-difference}中三支路电流分别是 $(V_n-V_{n-1})/r,(V_n-V_{n+1})/r,V_n/R$。它们的代数和为零；对从源点向右衰减的电势，左支路参考电流为负，表示实际从左流入。递推关系因而同时保留了方向和局部电荷守恒。

整理为 $V_{n+1}-(2+r/R)V_n+V_{n-1}=0$。设 $V_n=C\lambda^n$，得到特征方程
\[
 \lambda^2-(2+r/R)\lambda+1=0.
\]
两根之积为 $1$。在半无限网络中远处电势有界并趋向公共导线电势，排除大于 $1$ 的根；保留 $0<\lambda<1$ 的衰减根。有限网络中两根一般都需要，系数由两端边界确定，不能预先把增长根系数置零。

\fig{fig:c4:cut}{两条有阻导轨的切口}{\begin{tikzpicture}[scale=.88]
\draw(0.0,0)to[R,l={$r$}](0.0,2);\draw(0.0,2)to[R,l={$r$}](1.7,2);\draw(0.0,0)to[R,l_={$r$}](1.7,0);\draw(1.7,0)to[R,l={$r$}](1.7,2);\draw(1.7,2)to[R,l={$r$}](3.4,2);\draw(1.7,0)to[R,l_={$r$}](3.4,0);\draw(3.4,0)to[R,l={$r$}](3.4,2);\draw(3.4,2)to[R,l={$r$}](5.1,2);\draw(3.4,0)to[R,l_={$r$}](5.1,0);\draw(5.1,0)to[R,l={$r$}](5.1,2);\draw[dashed](5.1,0)--(6.1,0)(5.1,2)--(6.1,2);\node[left]at(0,2){$A$};\node[left]at(0,0){$B$};\draw[guide](2.2,-.6)--(2.2,2.7);\node at(3.8,3){保留右部两端口};
\end{tikzpicture}}
图\ref{fig:c4:cut}的两条导轨都含电阻。它的差电势递推应从上下两个节点的守恒相减得到；不能把下轨各点当作同一零电势节点。

双导轨网络中上、下每段均为 $r$，横档为 $R$。反向注入两导轨时，取差模电压 $V_n=V_n^{\text{上}}-V_n^{\text{下}}$，两导轨电势为 $\pm V_n/2$。图\ref{fig:c4:cut}中无源上节点的守恒式是
\[
 \frac{V_n-V_{n-1}}{2r}+\frac{V_n-V_{n+1}}{2r}+\frac{V_n}{R}=0.
\]
乘以 $2r$ 并把相邻列项移到同一边，得到
\[
 V_{n+1}-(2+2r/R)V_n+V_{n-1}=0.
\]
与单导轨多出的因子 $2$ 来自两根导轨都产生电势降。若误把理想回流线当成有阻导轨，会把衰减根和所有端口电阻一起算错。


有限网络的通解应写成 $A\lambda^n+B\lambda^{-n}$，再用真实远端边界确定两系数。远端移到第 $N$ 列后，常见开路边界使 $B/A$ 含 $\lambda^{2N}$ 的因子，所以在固定的近端位置，增长解的贡献才随 $N\to\infty$ 消失。直接删去增长根会丢失有限网络的边界影响。


\subsection{边界条件决定归一化系数}

对于双向无限双导轨，从第零横档上端注入 $I$、下端流出 $I$，左右对称给出 $V_n=V_0\lambda^{|n|}$，其中
\[
 \lambda=1+\frac rR-\sqrt{\frac{r^2}{R^2}+\frac{2r}{R}}.
\]
\fig{c4:source-normalization}{双向双导轨的源节点；上下势差为 $V_0$，每条上轨压降只有差电势变化的一半}{\begin{tikzpicture}

\draw(-3,2)to[R,l={$r$}](0,2)to[R,l={$r$}](3,2);
\draw(-3,-2)to[R,l_={$r$}](0,-2)to[R,l_={$r$}](3,-2);
\foreach \x in {-3,0,3}{\draw(\x,2)to[R,l={$R$}](\x,-2);}
\draw[dashed](-4,2)--(-3,2)(3,2)--(4,2)(-4,-2)--(-3,-2)(3,-2)--(4,-2);
\fill(0,2)circle(1.5pt)node[above right]{$V_0/2$};
\fill(0,-2)circle(1.5pt)node[below right]{$-V_0/2$};
\node[above]at(-3,2){$V_1/2$};\node[above]at(3,2){$V_1/2$};
\draw[flow](0,3.5)--(0,2.3)node[pos=0,above]{$I$};
\draw[flow](0,-2.3)--(0,-3.5)node[below]{$I$};
\draw[flow](.3,2.8)--(1.6,2.8)node[midway,above]{$K$};
\draw[flow](-.3,2.8)--(-1.6,2.8)node[midway,above]{$K$};
\draw[flow](.65,.8)--(.65,-.8)node[midway,right]{$V_0/R$};

\end{tikzpicture}}
图\ref{c4:source-normalization}中 $K=(V_0-V_1)/(2r)$，源节点满足 $I=2K+V_0/R$。利用左右对称把相邻列的差电势都写为 $V_1$。上轨从零列到相邻列的实际压降为 $(V_0-V_1)/2$，故不能直接用差电势除以 $r$；横档两端却正好相差 $V_0$。这正是归一化方程中两个不同系数的来源。

在上端源节点，两侧导轨分别流出 $(V_0-V_1)/(2r)$，横档流出 $V_0/R$。因此
\[
 I=\frac{V_0(1-\lambda)}r+\frac{V_0}R.
\]
将衰减根中的 $1-\lambda=-r/R+\sqrt{r^2+2rR}/R$ 代入，
\[
 \frac{I}{V_0}=\frac1R-\frac1R+
 \frac{\sqrt{r^2+2rR}}{rR}
 =\frac{\sqrt{r^2+2rR}}{rR},\qquad
 \frac{V_0}{I}=\frac{rR}{\sqrt{r^2+2rR}}.
\]
系数来自源节点的守恒，不是只解齐次递推就能确定。求出后还应检查所有横档电流之和等于总注入电流，确认没有漏掉左右两侧。


\section{虚拟电流与平移网络}
\subsection{叠加抵消非目标节点的注入}

线性网络可先解一组便于对称分析的外部电流，再把几组解相加。对于无限平面方格，可先考虑从 $A$ 注入 $I$、在很远边界分散流出的极限；旋转对称使源点四条边各流出 $I/4$。再从远边界注入、在相邻点 $B$ 流出 $I$，该点四条相邻边各流入 $I/4$。相加后边界电流抵消，端口为 $A,B$，直接边电流为 $I/2$，电势差为 $Ir/2$。

这里“流向无穷远”是局部电流的截断极限，而不是把无穷远定义成一个具有有限接地电阻的真实节点。在二维无限方格中，单源对极远边界的绝对电势差会发散，局部相邻电势差和配对源汇问题仍可有有限极限。推导只用这些局部电流差，不能据此声称无限方格源点对无穷远有某个有限电阻。

\fig{c4:lattice-auxiliary}{方格网络两辅助源汇的局部电流；所有边电阻为 $r$，四周继续延伸}{\begin{tikzpicture}

\foreach \shift/\mode in {0/1,8/2}{\begin{scope}[xshift=\shift cm,scale=.72]
\foreach \y in {-2,0,2}{\draw[gray](-2,\y)--(4,\y);\draw[dashed](-2.8,\y)--(-2,\y)(4,\y)--(4.8,\y);}
\foreach \x in {-2,0,2,4}{\draw[gray](\x,-2)--(\x,2);\draw[dashed](\x,-2.8)--(\x,-2)(\x,2)--(\x,2.8);}
\fill(0,0)circle(2pt)node[below left]{$A$};\fill(2,0)circle(2pt)node[below right]{$B$};
\ifnum\mode=1
\foreach \x/\y in {1/0,-1/0,0/1,0/-1}{\draw[flow,blue](0,0)--(\x,\y);}
\node at(1,3.6){$A$ 注入，四边各 $I/4$};
\else
\foreach \x/\y in {1/0,3/0,2/1,2/-1}{\draw[flow,blue](\x,\y)--(2,0);}
\node at(1,3.6){$B$ 流出，四边各 $I/4$};
\fi
\end{scope}}
\node at(4,-3){两图 $AB$ 电流同向，叠加后为 $I/2$；远边界电流相消};

\end{tikzpicture}}
图\ref{c4:lattice-auxiliary}左图只在 $A$ 附近使用四重旋转，右图只在 $B$ 附近使用四重旋转。最终源汇同时存在时，$A$ 的四条边已不等价；但同一条 $AB$ 边的两个贡献可以相加。端口电流仍是 $I$，于是直接边的压降给出整个网络两端的电势差 $Ir/2$。

\subsection{平移与局部旋转对称的条件}

无限蜂窝、三角和简单立方晶格也可作类似构造。保持边界截断方式一致，并在极限中恢复相应的平移、旋转对称。若每节点有 $z$ 条等价的电阻 $r$，从一个节点注入的电流在相邻边均分为 $I/z$；相邻源汇叠加后直接边为 $2I/z$，故相邻端口电阻为 $2r/z$。

该结论依赖所讨论的均匀欧氏周期晶格及相容边界极限，不能只看节点度数就推广到任意无限图。树形网络、存在特殊边界的网络或阻值不同的各向异性网络，可能不满足同样的电流分配。每次使用均分，都应在题设图中指出互换这些支路而保留其余条件的对称操作。


\section{自相似平面网络与三端结构}
\subsection{缩放的是几何，也是电阻}

如果均匀导线电阻与长度成正比，几何缩小为 $\eta$ 倍且截面、电阻率不变，整个相似子网络的端口电阻也缩小为 $\eta$ 倍。递推中写 $\eta X$ 的依据是每条导线电阻同时乘 $\eta$，不能只凭图形相似忽略阻值变化。

三角形或正方形嵌套结构常保留三个或四个连接点。把内部子图替换成一个两端电阻前，必须证明其他接点没有净外部电流，或使用完整的星形等效。对称轴上的节点虽然可能同电势，但拆开后各子网络是否仍保持该电势，需要检查各自的对称条件。保留三端响应的星形等效经常比勉强断开一个接点更可靠。


对具有三重对称、无源互易的三端子网络，可以选三条相等星臂为递推参数 $s$。星心电势由三个支路电流和为零确定：
\[
 V_O=\frac{V_1+V_2+V_3}{3},\qquad I_i=\frac{V_i-V_O}{s}.
\]
这组关系同时规定了三个端点的响应，比单次第三端悬空测得的两端电阻 $2s$ 保留更多接线信息。
\fig{c4:three-port-parallel}{两组三端等臂星形接在相同的三个端点；两个中心都处于三端电势的平均值}{\begin{tikzpicture}

\coordinate(O)at(0,0);\coordinate(Q)at(6,0);
\foreach \y/\lab in {3/1,0/2,-3/3}{
\coordinate(P\lab)at(3,\y);
\draw(O)to[R,l={$s_1$}](P\lab)to[R,l={$s_2$}](Q);
\fill(P\lab)circle(1.5pt)node[above right]{$V_{\lab}$};}
\fill(O)circle(1.5pt)node[left]{$O_1:\bar V$};\fill(Q)circle(1.5pt)node[right]{$O_2:\bar V$};
\node at(3,-4){合并 $O_1,O_2$ 后，各端对应星臂为 $s_1\parallel s_2$};

\end{tikzpicture}}
图\ref{c4:three-port-parallel}中两个星心等势可由上述电流方程分别推出，故可以连接并合并，再把同端两臂并联。该结论要求每个子网络的三条星臂各自相等；若三端不再对称，就须重新保留完整的三端响应，不能按同一标量递推。

\subsection{有限截断核验递推根}

得到自相似代数方程后，应列出几级有限截断值，检查其正性和趋近方向。新增内部正电阻通道使端口电阻不增，若计算序列反向变化，需检查是否误改了末端边界。若出现多个正代数根，有限截断与能量条件可以确定可实现分支，不能只选择看起来较简单的根。

电阻网络中的空间衰减还帮助判断截断误差：衰减根越小，远处边界对输入端的影响越弱。它与时间暂态的指数衰减形式相似，但自变量分别是节点编号与时间，物理约束不同。后续暂态讲义将用初始条件代替这里的空间边界条件。


\section{多导轨网络的对称分解}
\subsection{同向与反向电流分开处理}

上下两导轨完全相同。对任意外部注入，可把电流分成两轨同向的一部分和反向的一部分。相同方向、相同大小的激励使两轨电势相等，横档无电流；反向激励使两轨电势互为相反数，可以用差电势递推。最后将两部分电势和电流相加，恢复真实边界。

例如上轨零号节点流入 $I$、上轨 $n$ 号流出 $I$，可拆成：零号上下各流入 $I/2$、$n$ 号上下各流出 $I/2$ 的同向部分；加上零号上流入下流出 $I/2$、$n$ 号上流出下流入 $I/2$ 的反向部分。前者在两轨各产生 $n$ 段、每段电流 $I/2$ 的压降；后者用一个横档源解及其平移负解表示。这样既保留有阻回流，又避免把它误当成理想公共线。

对接地单导轨，常量模被横向接地电阻消除，远处差模指数衰减，端口距离足够大后电阻趋于饱和。对双导轨，平均电势的常量与线性模仍存在，异轨差模却衰减，因此沿同一条轨道相距很远的端口电阻还会含随距离增长的项。把两种模式混在一起，容易误判所有无限网络电阻都趋于常数。

\subsection{三条导轨的平均电势与偏差}

三条相同导轨每段电阻 $r$，同一列的三个节点两两以 $r$ 相连。令该列平均电势为 $\bar V_n$，各轨偏差为 $\delta_{a,n}=V_{a,n}-\bar V_n$，有 $\delta_{1,n}+\delta_{2,n}+\delta_{3,n}=0$。列内流出某节点的横向电流为
\[
 \frac{V_a-V_b}r+\frac{V_a-V_c}r=\frac{3\delta_a}{r}.
\]
\fig{c4:three-rail-column}{三导轨的同一列构成完整三角形；图中只展开上轨节点的左右邻边}{\begin{tikzpicture}

\coordinate(A)at(0,2);\coordinate(B)at(-2,-1);\coordinate(C)at(2,-1);
\draw(A)to[R,l={$r$}](B)to[R,l_={$r$}](C)to[R,l={$r$}](A);
\draw(A)to[R,l={$r$}](-3,4)node[left]{$V_{1,n-1}$};
\draw(A)to[R,l={$r$}](3,4)node[right]{$V_{1,n+1}$};
\fill(A)circle(1.5pt);\draw[guide](0,2.1)--(0,4.8);\node[above]at(0,4.8){$V_{1,n}=\bar V_n+\delta_{1,n}$};
\fill(B)circle(1.5pt)node[below left]{$V_{2,n}$};\fill(C)circle(1.5pt)node[below right]{$V_{3,n}$};
\node at(0,-2.3){本列横向流出和 $=(3V_{1,n}-V_{1,n}-V_{2,n}-V_{3,n})/r$};

\end{tikzpicture}}
图\ref{c4:three-rail-column}中，同列另外两点的电势和为 $3\bar V_n-V_{1,n}$，故上轨横向流出总量为 $3\delta_{1,n}/r$。左右纵向支路另外贡献 $(2\delta_{1,n}-\delta_{1,n-1}-\delta_{1,n+1})/r$，两部分合起来才得到系数 $5$。平均电势方程则没有这一横向项。


平均部分没有横向电流，相当于三条并联纵向通道。偏差部分的节点方程则为
\[
 5\delta_{a,n}-\delta_{a,n-1}-\delta_{a,n+1}=rJ_{a,n},
\]
其中 $J_{a,n}$ 是从该节点注入的电流减去同列平均注入。对某一轨的偏差分量，令零列偏差注入为单位电流，并使其余分量的偏差注入和为零。把该轨的电势偏差除以此电流，记响应电阻为 $G_n$。无源区先代入 $G_n=C\lambda^n$，除以 $C\lambda^{n-1}$ 得到
\[
 5\lambda-1-\lambda^2=0,\qquad
 \lambda_{\pm}=\frac{5\pm\sqrt{25-4}}2.
\]
两根互为倒数；远处偏差趋零，故右半轴只保留 $\lambda=(5-\sqrt{21})/2<1$ 的解，左半轴由反射得到 $G_n=G_0\lambda^{|n|}$。源点不能使用齐次方程，而应保留注入项：
\[
 5G_0-G_{-1}-G_1=r
 \quad\Longrightarrow\quad(5-2\lambda)G_0=r.
\]
由于 $5-2\lambda=\sqrt{21}$，幅度为 $G_0=r/\sqrt{21}$，所以
\[
 G_n=\frac{r}{\sqrt{21}}\lambda^{|n|}.
\]
对不同列的若干偏差电流乘以对应平移响应，再叠加平均电势，就能求任意端口。这一分解直接来自三轨置换对称和节点守恒，不需要把原网络错误折叠成一条导线。
